% ======================================================================
% Introduction.
% Research-paper funnel: field stake -> the bottleneck -> what prior work
% solved -> the capability still missing and what it costs -> this study.
% Implementation detail (identifiers, fingerprint protocol, admission rules)
% belongs to Methods and the Supplementary Information, not here.
% ======================================================================
\section{Introduction}

Materials discovery has become an exercise in reusing other people's
calculations. Public repositories now hold millions of computed crystals
\cite{jain2013commentary,saal2013materials,curtarolo2012aflow,schmidt2022dataset},
generative models propose new ones directly
\cite{xie2022crystal,zeni2025generative,merchant2023scaling}, and the universal
interatomic potentials used to filter those proposals are trained on whatever
the community has already published
\cite{chen2022universal,deng2023chgnet,batatia2023foundation,yang2024mattersim,wood2025uma}.
None of these activities consumes a single repository. Each draws on the union
of them, and each therefore rests on an assumption about that union: that it is
known what it contains, and that the numbers inside it mean the same thing.

Neither assumption has been tested, because the union has never been formed. The
repositories describe overlapping sets of crystals, but they were built
independently and share no key at the level of the object they describe. The same
arrangement of atoms carries a different accession number in every database that
holds it, and the one property that is genuinely shared, the chemical formula,
cannot serve as a key, because polymorphs share a formula and differ in almost
everything a screening campaign cares about. As a result three quantities that
bear directly on how materials data is used remain unmeasured. How much of the
corpus is the same crystal counted repeatedly. Whether two groups that computed
the same quantity for the same crystal obtained the same answer. And how the
properties a discovery campaign actually needs are distributed over the structures
that exist.

Forming the union once would in any case not be enough. Every repository
publishes new releases, and new repositories appear faster than any is retired,
so a union assembled in a single pass is out of date before it is finished and
has to be rebuilt from nothing to incorporate what arrived while it was being
made. A corpus that is to stay current has to do something a merge does not: it
has to recognise material it has never seen, and fold it in without discarding
the ability to recognise that it has already acquired.

Each of these gaps has already produced visible difficulty. Counts of newly
discovered stable materials cannot be reconciled between catalogues, and the
scale of recent generative releases has drawn precisely this criticism
\cite{cheetham2024artificial,leeman2024challenges,smit2026dataonly}. Benchmarks
that train on one repository and test on another inherit an unquantified overlap
term, which is why the Matbench Discovery protocol had to define its splits over
structures rather than records \cite{riebesell2025framework}. And a value that
two independent groups computed differently is, in the present corpus,
indistinguishable from a value that only one group ever computed, so no reader
can tell corroboration from coincidence.

Substantial progress has been made on the parts of this problem that can be
solved one layer at a time. The OPTIMADE specification gave more than twenty
providers a common query interface, removing the burden of a separate account
and format for each \cite{andersen2021optimade,evans2024developments}. Structural
fingerprints made deduplication tractable at scale, and the LeMaterial initiative
used one to publish a deduplicated union of several large repositories
\cite{ramlaoui2025lemat}. What neither delivers is the layer where the difficulty
actually resides. Knowing that two records describe the same crystal does not
attach what each database says about it, and attaching those statements is
useless unless they stay apart afterwards. Two repositories publishing a column
labelled band gap are frequently not reporting the same physical quantity: a
semilocal, a hybrid and a measured value differ by more than the differences a
screening campaign is trying to resolve. A union built by matching column names
would average them, and would destroy the very disagreement that makes the
corpus worth auditing.

Here we assemble that layer and use it to characterise the corpus. We resolve
identity across 54 databases at the level of the crystal structure, and retain for
every one of the resulting property values the database, the native record and the
calculation or measurement that produced it, so that several databases speaking
about one crystal remain separately audible. The construction answers the second
half of the problem as well, because admission and recognition are made the same
operation. A record that matches nothing is not an error to be logged but the
corpus's next member, and the moment it is admitted its structural fingerprint
joins the criterion against which every subsequent record is judged. The database
therefore evolves as it ingests: each round of assimilation enlarges what the next
round can recognise, so a newly released version or an entirely new repository is
folded into a substrate that has grown since the last one, rather than matched
against a frozen reference. We find that redundancy is large and
strongly concentrated, that independent databases reporting the same quantity
disagree by margins exceeding the thresholds on which materials selection
routinely turns, that a substantial part of the apparent agreement between them
is shared provenance rather than independent confirmation, and that coverage of
the property space is governed by what a property costs to compute rather than by
what it is worth knowing. We also find that the self-extension is
substantial rather than nominal: a sixth of all cross-database links exist only
because earlier assimilation had already created the structure they point to, and
for repositories whose chemistry the starting pool did not cover the figure
approaches unity. The same relation reappears at the next level up, where the
learned completion of the property layer is itself staged, each stage widening
the domain on which the next is defined. The instrument and the corpus it
produces are described in Methods and in the Supplementary Information; what
follows is what the corpus turns out to contain.
